% Source: openai/math@adc7f1241b42e322a6451854ab7e4b4c146bf78a,
% 07-comparators.tex, lines 255--298 and 332--354.
\section{An auxiliary label sequence for the selected state}

\begin{theorem}[Auxiliary mass and entropy along one label sequence]
    \label{thm:compressed_typical_label_sequence}
    \lean{FiniteProduct.exists_label_sequence_compressedTypicalSite_mass_entropy_asymptotic}
    \leanok
    \uses{thm:compressed_typical_site_norm,thm:high_label_pure_sequence,
        thm:compressed_site_label_mass_coordinates}
    Let $V$ be a finite set of sites with finite-dimensional coordinate
    spaces, let $X\subseteq V$, and let $\Omega\in\mathcal H_X\otimes
    \mathcal H_{X^c}$. Select eigenvalue indices $E$ of its actual
    marginal $\rho_X^\Omega$, with
    \begin{align*}
        d&=|E|,
        &z&=\sum_{i\in E}p_i>0.
    \end{align*}
    In the fixed enumeration of $E$, let
    $J_E:\mathbb C^d\to\mathcal H_X$ have the selected orthonormal
    eigenvectors as columns. Put $\mathcal H_R=\mathbb C^d$ and define
    \begin{align*}
        \chi&=z^{-1/2}(J_E^*\otimes I)\Omega,
        &\psi(b,r)&=\chi(r,b).
    \end{align*}
    Thus $\psi\in\mathcal H_{X^c}\otimes\mathcal H_R$ is unit. Its
    auxiliary density is
    \begin{align*}
        \rho_R&=\operatorname{Tr}_{\mathcal H_{X^c}}
            |\psi\rangle\langle\psi|,
    \end{align*}
    where the trace is over the physical first factor.
    Place the selected coordinate of $\chi$ at the auxiliary site $R$
    and leave its physical coordinates at their sites in $X^c$; write
    $\chi_{\mathrm{site}}$ for this vector and
    $\Psi_k=\chi_{\mathrm{site}}^{\otimes k}$ for its literal tensor power.

    There is one sequence of irreducible labels $\lambda_k$ of
    $\mathfrak S_k$ such that, with $Q_k$ the auxiliary isotypic
    projection of label $\lambda_k$,
    \begin{align*}
        q_k:=\|Q_k\Psi_k\|^2
        &\ge\frac{1}{2(k+1)^{d^2}}
            \quad\text{for all sufficiently large }k,\\
        \log\dim\lambda_k-kS(\rho_R)&=o(k).
    \end{align*}
    No normalization of the original vector $\Omega$ is required.
    This is the selected-state conclusion of
    \cite[Section~7, lines~255--281]{OpenAI2026AreaLaw}.
\end{theorem}
\begin{proof}
    \leanok
    \uses{thm:compressed_typical_site_norm,thm:high_label_pure_sequence,
        thm:compressed_site_label_mass_coordinates}
    Positive selected mass normalizes the actual compressed vector.
    Regroup its factors so that the physical complement comes first,
    and apply the pure-tensor-power label-sequence theorem to $\psi$.
    The auxiliary projection-mass identity preserves the selected label
    and its squared mass when returning to the site description.
    The auxiliary dimension is $d$, giving the displayed polynomial.
    The same selection supplies the logarithmic-dimension asymptotic
    for the actual density $\rho_R$.
\end{proof}

\begin{theorem}[One selected sequence and a common-cutoff Bell vector]
    \label{thm:compressed_typical_label_sequence_common_bell_vector}
    \lean{FiniteProduct.exists_label_sequence_commonLabelCutoff_compressedTypical_bell_vector}
    \leanok
    \uses{thm:compressed_typical_label_sequence,
        thm:compressed_typical_common_bell_vector}
    Use the preceding construction of $\psi$, $\rho_R$ and $\Psi_k$,
    and assume that the original physical vector $\Omega$ is unit.
    Fix a finite list
    $\mathcal B=(B_1,\ldots,B_m)$ of physical regions disjoint from $X$,
    every two of which are comparable by inclusion. The list may be empty
    or contain repetitions.
    Let $\mathcal H_C=\mathbb C^E$, and let the actual selected Bell vector
    and its projection be
    \begin{align*}
        b_E&=\frac{1}{\sqrt d}\sum_{i\in E}u_i\otimes|i\rangle_C,
        &\Pi_E&=|b_E\rangle\langle b_E|.
    \end{align*}
    Here $u_i$ is the selected eigenvector of $\rho_X^\Omega$;
    the $C$ coordinate is indexed by $E$, whereas the $R$ coordinate
    uses its fixed enumeration by $\{0,\ldots,d-1\}$.
    Put
    $\mathrm{post}_k=b_E^{\otimes k}\otimes\Psi_k$.
    For each $B_j$, let $\widehat B_j$ be its unchanged physical sites
    in the augmented family $\{R\}\sqcup X^c$, and define
    \begin{align*}
        C_{j,k}&=\mathbf 1_{(-\infty,\,k(S_\Omega(B_j)/z+1)]}
            (F_{\widehat B_j,k}),
        &C_k&=\prod_{j=1}^m C_{j,k}.
    \end{align*}
    The central label observable $F_{\widehat B_j,k}$ has value
    $\log\dim\mu$ on the isotypic subspace of label $\mu$.
    The product follows the prescribed list order; for the empty list
    it is the identity.

    There is one sequence $\lambda_k$, with its auxiliary projection
    $Q_k$, such that
    \begin{align*}
        \log\dim\lambda_k-kS(\rho_R)&=o(k).
    \end{align*}
    Define the masses using this same sequence by
    \begin{align*}
        q_k&=\|Q_k\Psi_k\|^2,
        &r_k&=\|C_kQ_k\Psi_k\|^2.
    \end{align*}
    For all sufficiently large $k$, there is a unit vector
    $\varphi_k\in(\mathcal H_X\otimes\mathcal H_C)^{\otimes k}
    \otimes\mathcal H_{\{R\}\sqcup X^c}^{\otimes k}$ satisfying
    \begin{align*}
        (\Pi_E^{\otimes k}\otimes I)\varphi_k&=\varphi_k,
        &(I\otimes Q_k)\varphi_k&=\varphi_k,\\
        (I\otimes C_{j,k})\varphi_k&=\varphi_k
            &&(1\le j\le m),\\
        \bigl|\langle\varphi_k,(I\otimes Q_k)
            \mathrm{post}_k\rangle\bigr|^2&=r_k,\\
        r_k&\ge q_k/2,
        &r_k&\ge\frac{1}{4(k+1)^{d^2}}.
    \end{align*}
    The vector $\varphi_k$ is invariant under simultaneous permutations
    of all copies, using the canonical identification between copies of
    the product space and the product of the copy spaces. The eventual
    threshold may depend on the fixed state, selection and region list.
    This is the common-cutoff overlap construction in
    \cite[Section~7, lines~283--298 and 332--354]{OpenAI2026AreaLaw}.
\end{theorem}
\begin{proof}
    \leanok
    \uses{thm:compressed_typical_label_sequence,
        thm:compressed_typical_common_bell_vector}
    Apply the preceding theorem to the actual compressed selected state,
    obtaining its one label sequence, eventual mass bound and entropy
    asymptotic. Keep that sequence in the common-cutoff Bell-vector theorem
    for the prescribed region list. The exponentially small regional losses
    leave at least half of its auxiliary mass, hence the displayed
    quarter-polynomial bound. Normalizing the common projection and adjoining
    the actual selected Bell power gives the four range and invariance
    properties. Its squared overlap with the unnormalized projected
    $\mathrm{post}_k$ is exactly the retained mass $r_k$.
\end{proof}
